\section{Discussion}
\label{sec:discussion}

The evaluation results show that the proposed enhancement improves several aspects of networking service quality in dynamic high-latency environments. The throughput improvement in static scenarios suggests that the additional control logic does not degrade the original steady-state behavior. The reduction in recovery time and throughput deficit area indicates faster adaptation after link changes. The decrease in p95 queue delay shows that the enhanced system also improves tail delay behavior rather than merely increasing throughput. Finally, the long-duration results demonstrate improved robustness under repeated dynamic switching.

The main reason for these improvements is the division of responsibilities among the three layers. The GRU-assisted perception layer improves the timeliness of link-state estimation, but it does not directly control packet transmission. The state-machine decision layer turns noisy estimates and feedback signals into interpretable control semantics. The execution layer then applies bounded corrections and protects the minimum FEC redundancy needed for streaming recovery. This structure avoids the risk of a fully neural control loop while still benefiting from temporal learning.

There are also limitations. First, the evaluation is based on Mininet emulation with controlled link scenarios. Although this setting is useful for repeatable experiments, it cannot cover all dynamics of real satellite or non-terrestrial networks. Second, the GRU model is trained from the available experimental traces, and its generalization to unseen network patterns requires further validation. Third, the current implementation focuses on a single long-lived TCP flow through the proxy. Multi-flow interaction, fairness, encrypted transport compatibility, and deployment issues remain future work.

Despite these limitations, the results indicate that lightweight learning-assisted perception can be useful when it is combined with explicit control semantics and conservative execution constraints. For PEP-based systems, this design direction is more practical than directly replacing the control loop with an opaque model, especially in high-latency environments where unstable actions may persist for one or more RTTs before their consequences become visible.
